\section{Limitations}
The write-stage estimate is conditional on paired provider completion. Six source trajectories are requested, but two failure-branch responses terminate before assistant text exists; GET-only recovery returns no memory text. Receipt-level forensics distinguish these missing arms from downstream parsing failure, but cannot rule out a label-dependent provider completion process or identify the direction of any resulting selection. We therefore report 4/4 divergence only among complete pairs and carry only those four source-memory pairs into downstream analysis.

The downstream evidence is intentionally finite and fixed-evidence. The larger first-action distribution stress test is unresolved ($p=0.311$), so the action claim remains an existence witness rather than population-level separation. The terminal result comes from a same-support uniform replication after an initial non-pass, retaining the identical four sources, four future tasks, 0.15 practical-effect floor, and $p<0.05$ gate while increasing only replication depth. Although the 256-call replication passes that gate, 8/16 cells are zero and two cells contain 78.2\% of squared-effect mass. The result therefore supports heterogeneous conditional sensitivity on this frozen matrix, not a stable task-level direction or a population effect. Source-faithful live browser navigation remains unexecuted because the released Shopping endpoint and reset service are unreachable from the execution hosts; this is transport debt, not negative scientific evidence.

Cross-writer evidence separates the write mechanism from downstream transport. GLM-5.3 changes all four paired memories and title sets, but its terminal replication yields 0.140625, below the preregistered 0.15 floor despite $p=0.00012$, with two direction reversals among six cells nonzero under both writers. Thus the write-time state intervention has bounded support across two writer families, while writer-invariant downstream magnitude or direction is not established. Additional writers and task domains are needed before making broader claims.

The no-memory arm is a source-independent branch-location control, not an independent $4\times4\times3$ factorial: each future-task baseline is reused descriptively across source rows, so we add no global significance test. Likewise, the corruption-rate curve is a deterministic consequence analysis of measured conditional effects, not an empirical corruption sweep. Under the released ReasoningBank retriever, label-invariant task descriptions are embedded with top-$1$ retrieval at threshold 0.3 and only the selected entry is injected; nonselected corruption-mask bits are therefore inert. A genuine multi-memory corruption interaction would require a different retrieval substrate. Token-set Jaccard and the fixed operation-slot taxonomy remain coarse diagnostics rather than learned semantic measures.

\section{Conclusion}
We ask a narrow causal question: when terminal feedback chooses how an episode is written into external memory, can an erroneous label become durable agent state even though the underlying experience is unchanged? On byte-identical released trajectories, switching only the reward-conditioned reflection branch changes every complete paired memory. A stronger same-mode semantic-preserving wording perturbation does not absorb the effect, supporting a reward-conditioned write channel rather than generic prompt sensitivity alone.

That state difference can reach later behavior under matched future evidence. Paired memories change 7/12 aligned next actions, and a same-support $4\times4$ terminal replication after an initial non-pass yields mean absolute success-rate difference 0.15625 with permutation $p=0.00074$. The effect is heterogeneous, and a no-memory control rules out a uniform ``success good / failure bad'' interpretation. GLM-5.3 reproduces the upstream write divergence but misses the frozen downstream practical-effect floor, so the evidence is stronger for the existence of the write-time channel than for invariant downstream magnitude or direction.

The resulting implication is configuration-bound but operationally important. In the released ReasoningBank/WebArena Shopping mechanism, terminal evaluation is part of state update: its label selects what durable memory is written, and that state can alter later outcome distributions. For reward-conditioned external-memory agents of this form, evaluator reliability is therefore also a persistent-state reliability requirement. Broader claims about writers, domains, live interaction, or corruption processes require separate evidence rather than extrapolation from this finite controlled study.
